\subsection{拉普拉斯积分变换的应用}
\label{subsec:applications}
拉普拉斯变换有很多应用,这里从三个方面的例子来说明.
\subsubsection{计算级数和}
有时,拉普拉斯变换可以用来计算某些级数的和.以下面的例子说明.

如计算级数和
$$
    \sum_{n=1} ^{\infty}  \frac{1}{n^2},
$$
由前面例题
$$
\int_0^{\infty} t e^{-p t} d t=\frac{1}{p^2}, \quad \Re p>0
$$

将级数化为
$$
\begin{aligned}
\sum_{n= 1}^{\infty} \frac{1}{n^2} & =\sum_{n=1}^{\infty} \int_0^{\infty} t e^{-n t} d t \\
& =\int_0^{\infty} t\left[\sum_{n=1}^{\infty} e^{-n t}\right] d t=\int_0^{\infty} \frac{t}{e^t-1} d t
\end{aligned}
$$
这个积分的值是 $\pi^2/6$,下面用 Parseval 等式核对,不必查表.

在 $[-\pi,\pi]$ 上取 $f(x)=x$,则 $\ell=\pi$,
\begin{equation}
  c_n = \frac{1}{2\pi}\int_{-\pi}^{\pi} x\, e^{-\imath n x}\,\mathrm{d}x.
\end{equation}
$n\neq 0$ 时分部积分.$\int_{-\pi}^{\pi} e^{-\imath n x}\,\mathrm{d}x=0$,只剩边界项.两端 $e^{-\imath n\pi}=e^{\imath n\pi}=(-1)^n$,并且 $1/(-\imath)=\imath$,
\begin{equation}
  \left[\frac{x\,e^{-\imath n x}}{-\imath n}\right]_{-\pi}^{\pi}
  = \frac{2\pi\,\imath}{n}(-1)^n.
\end{equation}
再除以 $2\pi$,得到 $c_n = \imath (-1)^n/n$,$|c_n|^2=1/n^2$,而 $c_0=0$.Parseval 等式
\begin{equation}
  \sum_{n=-\infty}^{\infty}|c_n|^2
  = \frac{1}{2\pi}\int_{-\pi}^{\pi} x^2\,\mathrm{d}x.
\end{equation}
左边是 $2\sum_{n=1}^{\infty} 1/n^2$.右边是
\begin{equation}
  \frac{1}{2\pi}\cdot\frac{2\pi^3}{3} = \frac{\pi^2}{3}.
\end{equation}
因此 $\sum_{n=1}^{\infty} 1/n^2=\pi^2/6$.左边的积分是 $\pi^2/3$,再除以 $2$ 才得到 $\pi^2/6$.

\subsubsection{求解定积分}
如果 $\int_v^{\infty} \bar{f}(q) d q$ 存在, 且当$t\to 0$时, $|f(t)/t|$有界,则
\begin{equation}
    \int_{p}^{\infty} \bar{f}(q) dq \risingdotseq \frac{f(t)}{t} . 
\end{equation}
比如
$$
\frac{\sin \omega t}{t} \fallingdotseq \int_p^{\infty} \frac{\omega}{q^2+\omega^2} d q=\frac{\pi}{2}-\arctan \frac{p}{\omega}.
$$

当$p\to 0$时, 有
\begin{equation}
    \int_0^{\infty} \bar{f}(p) d p=\int_0^{\infty} \frac{f(t)}{t} d t
\end{equation}
一个例子如
$$
\int_0^{\infty} \frac{\sin t}{t} d t=\int_0^{\infty} \frac{1}{p^2+1} d p=\frac{\pi}{2}
$$
不仅如此, 有些积分无法用留数定理计算,如
$$
\int_0^{\infty} \frac{\cos a t-\cos b t}{t} d t \quad a>0, b>0
$$
$\mathcal{L}\{\cos at\}=p/(p^2+a^2)$.于是
\begin{equation}
  \int_0^{\infty}\frac{\cos(at)-\cos(bt)}{t}\,\mathrm{d}t
  = \int_0^{\infty}\left(\frac{p}{p^2+a^2}-\frac{p}{p^2+b^2}\right)\mathrm{d}p.
\end{equation}
不定积分是 $\half\ln(p^2+a^2)$.定积分
\begin{equation}
  \left[\half\ln\frac{p^2+a^2}{p^2+b^2}\right]_0^{\infty}
  = 0 - \half\ln\frac{a^2}{b^2}
  = \ln\frac{b}{a}.
\end{equation}
上限处分子分母都像 $p^2$,对数趋于 $\ln 1=0$.

\subsubsection{求解微分方程}
给一个例子,利用拉普拉斯变换求解简谐振子方程的解.
\begin{example}
利用拉普拉斯变换求解简谐振子方程的解.
\end{example}
\begin{solution}
    质量为$m$的质点在弹性系数为$k$的弹簧牵引下做简谐运动,满足方程为
    $$
m \frac{d^2 X(t)}{d t^2}+k X(t)=0
$$
初始条件取
$$
X(0)=X_0, \quad X^{\prime}(0)=0
$$
应用拉氏变换到该方程上得到
$$
m \mathcal{L}\left\{\frac{d^2 X}{d t^2}\right\}+k \mathcal{L}\{X(t)\}=0
$$
用$x(p)$表示未知变换$\mathcal{L}\{ X(t)\}$,于是根据导数定理\eqref{eq:lapl_trans_derivative}有
$$
m p^2 x(p)-m p X_0+k x(p)=0,
$$
化简为
$$
x(p)=X_0 \frac{p}{p^2+\omega_0^2}, \quad \omega_0^2 \equiv \frac{k}{m} .
$$
查表得
$$
X(t)=X_0 \cos \omega_0 t.
$$
初速度为零,解里没有正弦项.
\end{solution}

若初速度改为 $X'(0)=V_0$,二阶导数定理多一项 $-m V_0$,于是
\begin{equation}
    x(p)=X_0\frac{p}{p^2+\omega_0^2}+V_0\frac{1}{p^2+\omega_0^2},
\end{equation}
\begin{equation}
    X(t)=X_0\cos(\omega_0 t)+\frac{V_0}{\omega_0}\sin(\omega_0 t).
\end{equation}
正弦项的系数是 $V_0/\omega_0$,因为 $\mathcal{L}\{\sin(\omega_0 t)\}=\omega_0/(p^2+\omega_0^2)$.在 $t=0$ 对解求导,可以核对这个因子.